% sec-A083833.tex -- k = 5 section of the unified kp+1
% paper. Standalone fragment, \input by
% A083833-UNIFIE-paper.tex.
% Environments and notation inherited from sec-carry.tex; the
% guards below are no-ops once defined there.
\providecommand{\IsPal}{\mathrm{IsPal}}
\makeatletter
\@ifundefined{theorem}{\newtheorem{theorem}{Theorem}}{}
\@ifundefined{corollary}{\newtheorem{corollary}[theorem]{Corollary}}{}
\@ifundefined{remark}{\newtheorem{remark}[theorem]{Remark}}{}
\makeatother

\section{Multiplier $k = 5$: A083833 and its image A083834}
\label{sec:k5}

OEIS A083833~\cite{oeis} (Murthy and Srikanth, 2003)
lists the palindromes $p$ such that $q = 5p+1$ is again
a palindrome; A083834 is the image sequence of the
values $q$. On A083833, Michael S.~Branicky recorded
(June 13, 2021, tagged ``(Conjectures)''): ``All terms
contain only the digits $\{0,1,2,3\}$'' and ``For $d$
even, the only term with $d$ digits is $2^d$;
equivalently, all terms starting with $1$ have odd
length'' (in the entry's notation the caret $2^d$
denotes the $d$-fold concatenated repdigit $22\ldots2$).
On the image entry A083834 he recorded the analogous
pair, the alphabet $\{0,1,5,6\}$ and a repunit family;
that verbatim contains a parity slip discussed in
Remark~\ref{rem:parity-slip}.

\begin{theorem}[Branicky C1: alphabet of $p$]
\label{thm:k5-c1}
For all naturals $p, m, m'$ with $\IsPal(p,m)$,
$\IsPal(5p+1,m')$ and $m \ge 1$, every digit of $p$
satisfies $a_j \le 3$ for all $j < m$.
\end{theorem}

\begin{theorem}[Branicky C2: even length, full form]
\label{thm:k5-c2}
Let $m \ge 1$ be even and $\IsPal(p,m)$. Then $5p+1$ is
a palindrome if and only if $a_j = 2$ for all $j < m$:
the even-length members are exactly the repdigits
$p = 22\ldots2$.
\end{theorem}

A term starting with $1$ is not the repdigit, so by
Theorem~\ref{thm:k5-c2} it cannot have even length:
this is Branicky's ``equivalently'' clause, a
paper-side corollary of C2 rather than a separate
kernel declaration.

\begin{proof}[Proof of
Theorems~\ref{thm:k5-c1}--\ref{thm:k5-c2} and
Corollary~\ref{cor:k5-image}]
The shared carry-chain device (Section~\ref{sec:carry})
instantiates as follows. The carry into position $j$ is
$c_j = \bigl\lfloor\bigl(5\,(p \bmod 10^j)+1\bigr)
/10^j\bigr\rfloor$, bounded by $4$, with $c_0 = 1$; the
output length satisfies $m' \in \{m, m+1\}$. If
$m' = m$, Lemma~\ref{lem:same} forces $m$ odd,
$a_0 = 1$, $c_1 = 0$, and the alternating classes
$a_j \in \{0,1\}$ at even $j$, $a_j \in \{2,3\}$ at odd
$j$. If $m' = m+1$, the swap-freeness collapses the
chain to the constant $1$ (Scheme II of
Lemma~\ref{lem:overflow}); the boundary equations give
$a_0 = 2$, and the parity propagation
$a_j \equiv a_{j-1} \pmod 2$ pins every digit to $2$:
$p = 22\ldots2$ and $q = 11\ldots1$ is a repunit.
Reading $b_j = (5a_j + c_j) \bmod 10$ off the two
branches gives the image alphabet $\{1,6\}$ at even $j$,
$\{0,5\}$ at odd $j$, and the repunit $1\ldots1$ at
overflow.
Theorem~\ref{thm:k5-c1} is the union of the two digit
classes; if $m$ is even the same-length alternative is
excluded, which is the forward direction of
Theorem~\ref{thm:k5-c2}; the converse direction is the
certified identity
$5\cdot 2R_m + 1 = R_{m+1}$ packaged as a literal iff
(Table~\ref{tab:lean}). The $29$ listed terms fit:
their even-length members are exactly $22$, $2222$,
$222222$, $22222222$.
\end{proof}

\begin{remark}[parity slip in the A083834 verbatim]
\label{rem:parity-slip}
On A083834 the comment reads ``For $d$ \emph{odd}, the
only term with $d$ digits is $1^d$; equivalently, all
terms starting with $6$ have odd length''. The ``$d$
odd'' is a slip: the listed odd-length terms $606$ and
$656$ are not repunits, so the literal reading is
refuted by the entry's own data, whereas the
``equivalently'' clause is unambiguous. The corrected
statement is that even-length terms are repunits, and
that is the statement the kernel proves. On A083833
the analogous comment had been recorded correctly with
``$d$ even''.
\end{remark}

\begin{corollary}[image alphabet: A083834]
\label{cor:k5-image}
Under the same hypotheses, every digit of $q = 5p+1$
lies in $\{0,1,5,6\}$; if $m'$ is even then every digit
of $q$ equals $1$, i.e.\ $q$ is the repunit
$11\ldots1$; and a term of A083834 starting with $6$
has odd length.
\end{corollary}
